% Shared exercise chunk: l2-transduction
\begin{exercisechunk}
\item \textbf{Induction and transduction.}
\label[exercise]{ex:transduction}
An \emph{inductive} learner uses the training sample to produce a prediction
function for future inputs. A \emph{transductive} learner also receives the
particular unlabeled test inputs whose labels it must predict.
\begin{enumerate}[(a)]
\item Consider the original boxes protocol with one test input $X$, and allow
any known promise $f\in\cH\subseteq\{0,1\}^{[N]}$. A transductive learner
$\cB$ receives $(s,x,w)$ and returns a bit. Show that there is an inductive
learner $\cA$ of the form \cref{eq:randomized-algorithm} with exactly the same
risk as $\cB$ for every $f\in\cH$. The training sample and its selection
procedure are unchanged, and the learner's output function need not belong
to $\cH$.

\item Three boxes have nondecreasing binary labels. The training data reveal
$f(1)=0$ and $f(3)=1$. The test set consists of box 2 and one other box, with
the promise that their labels differ. Loss is the fraction of the two test
labels predicted incorrectly. Allow randomized learners in both cases; the
hidden labels are chosen before the learner's private random seed is drawn.
\begin{enumerate}[label=(\roman*),leftmargin=*,beginpenalty=10000]
\item Find the smallest worst-case expected loss when the learner must produce its
prediction function before seeing the test set.
\item Then find it when the learner
receives both test indices before predicting their labels.
\end{enumerate}
\end{enumerate}

\emph{Where does the gain come from?}
In part~(b), the promise links the test-set composition to the hidden labels,
so the unlabeled test indices themselves convey information about $f$.
This changes the original assumption that index selection does not use $f$.
Part~(a) concerns a single test input with unchanged training data; it does
not rule out gains from seeing a whole unlabeled test set under additional
promises such as the one in part~(b).
\end{exercisechunk}
